% Eigenständige Lesefassung zur Kompaktheit in metrischen Räumen.
% Die nummerierten Aussagen werden ausschließlich im Hauptband deklariert.
\chapter*{Wie Vollständigkeit und endliche Netze zusammenwirken}
\hypertarget{compactness.reading}{}
\addcontentsline{toc}{chapter}{Wie Vollständigkeit und endliche Netze zusammenwirken}

Eine beschränkte Folge muss nicht gegen einen Punkt streben. Selbst wenn
der Raum vollständig ist, muss sie nicht einmal eine konvergente Teilfolge
besitzen. Das zeigt die Menge \(\mathbb N\) mit dem diskreten Abstand:
Verschiedene Punkte haben Abstand \(1\), gleiche Punkte Abstand \(0\).
Der ganze Raum ist beschränkt, und jede Cauchy-Folge wird schließlich
konstant. Trotzdem hat die Folge \(0,1,2,\ldots\) keine konvergente
Teilfolge. Ihre verschiedenen Glieder bleiben immer im Abstand \(1\).

Vollständigkeit garantiert also nur dann einen Grenzwert, wenn die
Folgenglieder bereits im Sinne von Cauchy zusammenrücken. Eine zusätzliche
Bedingung muss dafür sorgen, dass man aus jeder Folge wenigstens eine
solche Teilfolge herausgreifen kann. Genau das leistet die
\emph{totale Beschränktheit}. Sie verlangt, dass der Raum bei jeder noch
so feinen Auflösung durch endlich viele Kugeln überdeckt werden kann.

Wir werden zeigen, dass daraus zusammen mit Vollständigkeit genau die
Folgenkompaktheit entsteht. Anschließend betrachten wir offene
Überdeckungen. Auch deren endliche Reduzierbarkeit beschreibt in
metrischen Räumen dieselbe Eigenschaft.

\section*{Kugeln im betrachteten Raum}

Sei \((X,d)\) ein metrischer Raum und \(K\subseteq X\). Auf \(K\)
verwenden wir den eingeschränkten Abstand \(d_K\). Es werden keine
Abstände verändert; lediglich die zulässigen Punkte gehören jetzt zu
\(K\). Für \(c\in K\) und \(r>0\) schreiben wir
\[
 B_K(c,r)=\{x\in K\mid d(x,c)<r\}
          =K\cap B_d(c,r).
\]
Wenn von Vollständigkeit oder Konvergenz \emph{in \(K\)} die Rede ist,
müssen auch die Grenzwerte zu \(K\) gehören. So konvergiert die Folge
\(1/(n+2)\) in \(\mathbb R\) gegen \(0\), besitzt aber keinen Grenzwert
im Teilraum \((0,1)\).

Eine Teilmenge \(U\subseteq K\) heißt \emph{offen in \(K\)}, wenn zu
jedem \(x\in U\) ein \(r>0\) mit \(B_K(x,r)\subseteq U\) existiert.
Diese Formulierung ist relativ zum betrachteten Raum. Beispielsweise ist
\([0,\tfrac12)\) offen in \([0,1]\), obwohl es nicht offen in
\(\mathbb R\) ist. Am Punkt \(0\) wird nur die Kugel innerhalb von
\([0,1]\) betrachtet.

Wir indexieren Folgen durch \(\mathbb N=\{0,1,2,\ldots\}\). Eine
Teilfolge entsteht durch streng wachsende Indizes
\(j_0<j_1<j_2<\cdots\). Insbesondere gilt \(j_k\geq k\).
Diese einfache Beobachtung wird bei den kleiner werdenden Radien
\(1/(n+1)\) mehrfach nützlich sein.

\section*{Endliche Netze bei jeder Auflösung}

Ein \emph{endliches \(\varepsilon\)-Netz} von \(K\), wobei
\(\varepsilon>0\), ist eine endliche Menge \(F\subseteq K\) mit
\[
 K=\bigcup_{c\in F}B_K(c,\varepsilon).
\]
Jeder Punkt von \(K\) liegt also in einem Abstand kleiner als
\(\varepsilon\) von mindestens einem der endlich vielen Netzpunkte.
Die Zentren gehören selbst zu \(K\). Man kann die Netzpunkte als
endlich viele Vertreter auffassen: Bei der vorgegebenen Genauigkeit
liegt jeder Punkt nahe genug bei einem Vertreter.

\(K\) heißt \emph{total beschränkt}, wenn für jedes
\(\varepsilon>0\) ein solches endliches Netz existiert. Dabei darf
das Netz von \(\varepsilon\) abhängen. Die Zahl seiner Punkte braucht
keine gemeinsame obere Schranke für alle Genauigkeiten zu besitzen.
Dies sind die Definitionen
\FormulaRefAuto{MetricFiniteEpsilonNetDef}[def] und
\FormulaRefAuto{MetricTotalBoundednessDef}[def] des Hauptbandes.

Für eine endliche Menge \(K\) kann man stets \(F=K\) nehmen.
Auch die leere Menge ist total beschränkt: Das leere Netz genügt, denn
die Vereinigung einer leeren Familie ist leer. Diese Konvention
erspart keine Beweise für nichtleere Räume, sorgt aber dafür, dass die
späteren Sätze ohne unnötige Ausnahmen gelten.

Die Forderung ist stärker als gewöhnliche Beschränktheit. Ein
beschränkter Raum passt in eine einzige hinreichend große Kugel.
Ein total beschränkter Raum lässt sich zusätzlich mit endlich vielen
\emph{beliebig kleinen} Kugeln erfassen. Beim diskreten Raum
\(\mathbb N\) ist jede Kugel vom Radius \(1/2\) eine Einermenge.
Endlich viele solcher Kugeln können nicht alle natürlichen Zahlen
enthalten. Dieser Raum ist deshalb nicht total beschränkt.

Dass totale Beschränktheit tatsächlich gewöhnliche Beschränktheit
nach sich zieht, lässt sich unmittelbar nachrechnen. Sei \(K\ne
\varnothing\), und sei \(F\) ein endliches \(1\)-Netz. Dann ist
\(F\) nichtleer. Wähle \(c_0\in F\) und setze
\[
 M=\max\{d(c,c_0)\mid c\in F\}.
\]
Zu jedem \(x\in K\) gibt es ein \(c\in F\) mit \(d(x,c)<1\).
Die Dreiecksungleichung liefert
\[
 d(x,c_0)\leq d(x,c)+d(c,c_0)<1+M.
\]
Also liegt \(K\) in einer Kugel vom Radius \(1+M\). Für
\(K=\varnothing\) gilt die Beschränktheitsbedingung ohnehin.
Dies beweist \FormulaRefAuto{MetricTotalBoundedSetBounded}[thm].

\section*{Wenn ein endliches Netz unmöglich ist}

Fehlt totale Beschränktheit, so gibt es ein \(\varepsilon_0>0\),
für das kein endliches \(\varepsilon_0\)-Netz existiert. Wir können
dann eine Folge auswählen, deren verschiedene Glieder stets einen
Abstand von mindestens \(\varepsilon_0\) haben.

Der Raum \(K\) ist in diesem Fall nichtleer. Wähle zunächst
\(x_0\in K\). Sind \(x_0,\ldots,x_n\) bereits gewählt, so
überdecken die Kugeln \(B_K(x_i,\varepsilon_0)\) mit \(i\leq n\)
den Raum nicht. Deshalb kann der nächste Punkt außerhalb ihrer
Vereinigung gewählt werden:
\[
 x_{n+1}\in K\setminus
       \bigcup_{i=0}^{n}B_K(x_i,\varepsilon_0).
\]
Für alle verschiedenen \(i,j\) folgt daraus
\[
 d(x_i,x_j)\geq\varepsilon_0.
\]

Hier steckt eine Auswahlstelle: Aus jeder der auftretenden nichtleeren
Restmengen muss ein Punkt bestimmt werden. Das Auswahlprinzip liefert
eine Auswahlfunktion auf den nichtleeren Teilmengen von \(K\).
Mit dieser festen Funktion definiert der Rekursionssatz die Folge.
Es genügt also nicht, nur für jeden einzelnen Schritt die Existenz
eines nächsten Punktes zu bemerken; Auswahl und Rekursion verbinden
die Schritte zu einer Folge.

Keine Teilfolge dieser Folge ist Cauchy. Zu zwei verschiedenen
Teilfolgenindizes bleibt der Abstand mindestens
\(\varepsilon_0\), gleichgültig, wie weit hinten sie liegen.
Damit kann auch keine Teilfolge konvergieren, denn jede konvergente
Folge ist Cauchy. Wir haben bereits zwei Folgerungen gewonnen:
\[
 \begin{gathered}
 \text{Jede Folge in \(K\) besitzt eine Cauchy-Teilfolge}
       \ \Longrightarrow\ K\text{ ist total beschränkt},\\[.4em]
 K\text{ ist folgenkompakt}
       \ \Longrightarrow\ K\text{ ist total beschränkt}.
 \end{gathered}
\]
Die erste Folgerung wird gleich zu einer Äquivalenz ergänzt.

\section*{Aus einer Folge eine Cauchy-Teilfolge gewinnen}

Sei jetzt \(K\) total beschränkt, und sei \((x_n)_{n\in\mathbb N}\)
eine Folge in \(K\). Für \(K=\varnothing\) gibt es keine solche
Folge; die Aussage über alle Folgen gilt dort ohne weiteren Nachweis.
Im Folgenden ist also \(K\ne\varnothing\).

Wir überdecken \(K\) zunächst durch endlich viele Kugeln vom
Radius \(1\). Wenigstens eine dieser Kugeln enthält \(x_n\) für
unendlich viele Indizes \(n\). Andernfalls wäre die Vereinigung
der endlich vielen endlichen Indexmengen endlich und könnte nicht
ganz \(\mathbb N\) sein. Nenne die unendliche Indexmenge dieser
Kugel \(I_0\).

Im nächsten Schritt überdecken wir \(K\) durch endlich viele Kugeln
vom Radius \(1/2\). Wenigstens eine enthält \(x_n\) für unendlich
viele Indizes \(n\in I_0\). Die zugehörige Indexmenge sei
\(I_1\subseteq I_0\). So fahren wir mit den Radien
\(1/3,1/4,\ldots\) fort. Wir erhalten unendliche Mengen mit
\[
 I_0\supseteq I_1\supseteq I_2\supseteq\cdots
\]
und für jedes \(k\) ein Zentrum \(c_k\in K\), so dass
\[
 n\in I_k\quad\Longrightarrow\quad
              d(x_n,c_k)<\frac1{k+1}.
\]
Für je zwei Indizes \(p,q\in I_k\) gilt daher
\[
 d(x_p,x_q)
 \leq d(x_p,c_k)+d(c_k,x_q)<\frac2{k+1}.
\]
Die Durchmesser der so ausgewählten Punktmengen werden beliebig
klein. Verschachtelt sind dabei die \emph{Indexmengen} \(I_k\).
Die verwendeten Kugeln müssen nicht ineinander enthalten sein.

Auch hier ist die unendliche Auswahl zu begründen. Für jede Stufe
und jede bisherige unendliche Indexmenge existiert eine passende
Kugel mit unendlich vielen verbleibenden Indizes. Man kann mit
dem Auswahlprinzip eine Auswahlfunktion für diese zulässigen
Fortsetzungen festlegen und dann rekursiv die Mengen \(I_k\)
und Zentren \(c_k\) bilden. Das ist dieselbe Verbindung von
Auswahl und Rekursion wie im vorherigen Abschnitt.

Jetzt wählen wir einen einzelnen, hinreichend späten Index aus
jeder Stufe. Setze
\[
 j_0=\min I_0,\qquad
 j_{k+1}=\min\{n\in I_{k+1}\mid n>j_k\}.
\]
Die zweite Menge ist nichtleer: Eine unendliche Teilmenge von
\(\mathbb N\) kann nicht in einem endlichen Anfangsabschnitt
liegen. Ihre kleinsten Elemente sind eindeutig bestimmt, so dass
für diese letzte Rekursion keine weitere Auswahlfunktion nötig ist.
Die Indizes \(j_k\) wachsen streng.

Sei \(\varepsilon>0\). Wähle \(N\) so groß, dass
\(2/(N+1)<\varepsilon\). Für \(p,q\geq N\) gilt wegen der
Verschachtelung
\[
 j_p\in I_p\subseteq I_N,
 \qquad j_q\in I_q\subseteq I_N.
\]
Die eben bewiesene Abstandsabschätzung ergibt
\[
 d(x_{j_p},x_{j_q})<\frac2{N+1}<\varepsilon.
\]
Damit ist \((x_{j_k})\) eine Cauchy-Teilfolge. Zusammen mit dem
vorherigen Abschnitt ist bewiesen:
\[
 \boxed{\begin{gathered}
 K\text{ ist total beschränkt}\\[-.2em]
 \Longleftrightarrow\\[-.2em]
 \text{Jede Folge in \(K\) besitzt eine Cauchy-Teilfolge.}
 \end{gathered}}
\]
Dies ist das Kriterium
\FormulaRefAuto{MetricTotalBoundedCauchySubsequenceCriterion}[thm].
Zum ausführlichen Beweis:
\CompactnessProofLink{MetricTotalBoundedCauchySubsequenceCriterion}.

\section*{Vollständigkeit liefert den Grenzwert}

Die totale Beschränktheit hat uns eine Cauchy-Teilfolge geliefert.
Ob diese Teilfolge in \(K\) einen Grenzwert besitzt, ist eine andere
Frage. Genau dafür ist Vollständigkeit zuständig. Ist \((K,d_K)\)
vollständig und \(K\) total beschränkt, so besitzt jede Folge in
\(K\) eine Cauchy-Teilfolge, und diese konvergiert gegen einen
Punkt von \(K\). Somit ist \(K\) folgenkompakt.

Für die Umkehrung sei \(K\) folgenkompakt. Dass \(K\) total
beschränkt ist, folgt bereits aus der Konstruktion der getrennten
Folge. Zur Vollständigkeit betrachten wir eine Cauchy-Folge
\((x_n)\) in \(K\). Nach Folgenkompaktheit gibt es eine Teilfolge
\((x_{j_k})\), die gegen einen Punkt \(x\in K\) konvergiert.
Dann konvergiert sogar die ursprüngliche Folge gegen \(x\).

Um dies zu sehen, sei \(\varepsilon>0\). Die Cauchy-Eigenschaft
liefert ein \(N\) mit
\(d(x_m,x_n)<\varepsilon/2\) für \(m,n\geq N\). Wähle
danach ein \(k\) mit \(j_k\geq N\) und
\(d(x_{j_k},x)<\varepsilon/2\). Für jedes \(n\geq N\) gilt
\[
 d(x_n,x)\leq d(x_n,x_{j_k})+d(x_{j_k},x)<\varepsilon.
\]
Damit ist \((K,d_K)\) vollständig. Dies ist zugleich die
Beweisidee des bereits bekannten Satzes
\FormulaRefAuto{MetricSequentiallyCompactSetComplete}[thm].

Wir erhalten die Charakterisierung
\FormulaRefAuto{MetricSequentialCompactnessCharacterization}[thm]:
\[
 \boxed{\begin{gathered}
 K\text{ ist folgenkompakt}\\[-.2em]
 \Longleftrightarrow\\[-.2em]
 (K,d_K)\text{ ist vollständig und }K\text{ ist total beschränkt}.
 \end{gathered}}
\]
Die beiden Bedingungen erfüllen verschiedene Aufgaben. Totale
Beschränktheit sorgt für eine Cauchy-Teilfolge; Vollständigkeit
sorgt für deren Grenzwert innerhalb des Raums. Gewöhnliche
Beschränktheit kann die erste Aufgabe nicht übernehmen.

\section*{Drei Beispiele, die die Bedingungen trennen}

Die folgende Tabelle bezieht sich bei den Intervallen auf den
gewöhnlichen reellen Abstand und bei \(\mathbb N\) auf den
diskreten Abstand.

\begin{center}
\small
\begin{tabular}{@{}lcccc@{}}
\toprule
Raum & beschränkt & total beschränkt & vollständig & folgenkompakt\\
\midrule
\(\mathbb N\), diskret & ja & nein & ja & nein\\
\((0,1)\)             & ja & ja   & nein & nein\\
\([0,1]\)            & ja & ja   & ja & ja\\
\bottomrule
\end{tabular}
\end{center}

Beim diskreten Raum haben wir alle Einträge bereits begründet.
Die beiden Intervalle sind total beschränkt, weil sich ein
Intervall mit einem endlichen Gitter beliebig fein erfassen lässt.
Genauer: Wähle zu \(\varepsilon>0\) eine positive natürliche
Zahl \(m\) mit \(1/m<\varepsilon\). Die Mittelpunkte
\[
 c_i=\frac{2i+1}{2m}\qquad(0\leq i<m)
\]
liegen alle in \((0,1)\). Jeder Punkt von \([0,1]\) gehört zu
einem der Teilintervalle \([i/m,(i+1)/m]\) und hat von dessen
Mittelpunkt höchstens den Abstand \(1/(2m)<\varepsilon\).
Dieselben Zentren bilden daher sowohl für \([0,1]\) als auch
für \((0,1)\) ein endliches \(\varepsilon\)-Netz.

Das offene Intervall \((0,1)\) ist dennoch nicht vollständig.
Die Folge \(x_n=1/(n+2)\) liegt darin und ist Cauchy, weil sie
in \(\mathbb R\) gegen \(0\) konvergiert. Ein Grenzwert in
\((0,1)\) wäre zugleich ein reeller Grenzwert und müsste wegen
der Eindeutigkeit gleich \(0\) sein. Dieser Punkt fehlt.

Das abgeschlossene Intervall \([0,1]\) ist dagegen vollständig.
Jede Cauchy-Folge darin besitzt wegen der Vollständigkeit von
\(\mathbb R\) einen reellen Grenzwert \(x\). Wäre \(x<0\),
hätten alle Folgenglieder Abstand mindestens \(-x>0\) von \(x\);
wäre \(x>1\), hätten sie Abstand mindestens \(x-1>0\).
Beides widerspricht der Konvergenz. Also liegt \(x\in[0,1]\).
Der Hauptsatz liefert nun die Folgenkompaktheit des Intervalls.

\section*{Kompaktheit durch offene Überdeckungen}

Eine \emph{offene Überdeckung} von \(K\) ist eine Familie
\(\mathcal U\) von in \(K\) offenen Mengen mit
\[
 K=\bigcup_{U\in\mathcal U}U.
\]
Jeder Punkt wird von mindestens einer der Mengen erfasst.
Die Mengen dürfen unterschiedliche Formen und Größen haben,
und die Familie darf unendlich sein.

\(K\) heißt \emph{kompakt}, wenn zu jeder offenen Überdeckung
\(\mathcal U\) eine endliche Teilfamilie
\(\mathcal V\subseteq\mathcal U\) mit
\[
 K=\bigcup_{U\in\mathcal V}U
\]
existiert. Man darf dabei nur Mengen aus der gegebenen Überdeckung
behalten. Das ist die Definition
\FormulaRefAuto{MetricOpenCoverCompactnessDef}[def].

Für \(K=\varnothing\) genügt stets die leere Teilfamilie.
Die leere Menge ist also kompakt. Sie ist auch folgenkompakt
und vollständig, weil keine Folgen mit Werten in ihr existieren.
Damit gelten die folgenden Äquivalenzen auch für diesen Fall.

Endliche Netze sehen bereits wie endliche Überdeckungen aus.
Der Unterschied liegt in der Quantifizierung: Bei totaler
Beschränktheit suchen wir zu jedem Radius \emph{irgendeine}
endliche Kugelüberdeckung. Bei Kompaktheit muss \emph{jede}
vorgegebene offene Überdeckung eine endliche Teilüberdeckung
besitzen.

Am offenen Intervall wird dieser Unterschied sichtbar. Die Mengen
\[
 U_n=\bigl(1/(n+2),1\bigr)\qquad(n\in\mathbb N)
\]
sind offen in \((0,1)\) und überdecken es: Zu jedem \(x>0\)
gibt es ein \(n\) mit \(1/(n+2)<x\). Endlich viele dieser
Mengen genügen aber nicht. Ist \(N\) der größte ausgewählte
Index, ist ihre Vereinigung in \(U_N\) enthalten und verfehlt
beispielsweise \(1/(2(N+2))\). Also ist \((0,1)\) nicht
kompakt, obwohl es total beschränkt ist.

\section*{Eine gemeinsame Kugelgröße}

Sei \(K\ne\varnothing\) folgenkompakt, und sei
\(\mathcal U\) eine offene Überdeckung. Für jeden Punkt
\(x\in K\) gibt es zunächst eine passende Umgebung: Man wählt
ein \(U\in\mathcal U\) mit \(x\in U\), und wegen der
Offenheit passt eine hinreichend kleine Kugel um \(x\) ganz in
\(U\). Wir zeigen, dass eine positive Kugelgröße
\emph{gleichzeitig für alle Punkte} genügt:
\[
 \text{Es gibt }\delta>0\text{ mit }
 \forall x\in K\ \exists U\in\mathcal U:
                  B_K(x,\delta)\subseteq U.
\]
Eine solche Zahl nennen wir hier eine \emph{Lebesgue-Zahl}
in der Kugelformulierung. Es braucht nicht für alle Punkte
dasselbe Überdeckungselement \(U\) zu sein; gemeinsam ist nur
der Radius \(\delta\).

Angenommen, es gäbe keinen solchen Radius. Dann könnte man für
jedes \(n\) einen Punkt \(x_n\in K\) finden, für den die
Kugel \(B_K(x_n,1/(n+1))\) in keinem einzelnen
\(U\in\mathcal U\) enthalten ist. Das Auswahlprinzip liefert
aus diesen nichtleeren Mengen möglicher Gegenbeispiele eine
Folge \((x_n)\).

Wegen der Folgenkompaktheit hat sie eine Teilfolge
\(x_{j_k}\to x\) mit \(x\in K\). Wähle ein
\(U\in\mathcal U\) mit \(x\in U\) und danach \(r>0\)
mit \(B_K(x,r)\subseteq U\). Für hinreichend großes \(k\)
gelten gleichzeitig
\[
 d(x_{j_k},x)<r/2,
 \qquad \frac1{j_k+1}<r/2.
\]
Die zweite Ungleichung folgt aus \(j_k\geq k\). Für jeden
Punkt \(y\in B_K(x_{j_k},1/(j_k+1))\) gilt nun
\[
 d(y,x)\leq d(y,x_{j_k})+d(x_{j_k},x)
         <\frac1{j_k+1}+\frac r2<r.
\]
Somit liegt die ganze Kugel
\(B_K(x_{j_k},1/(j_k+1))\) in \(U\). Das widerspricht
der Wahl von \(x_{j_k}\). Die gemeinsame Größe \(\delta\)
existiert also.

Nun ist \(K\) als folgenkompakte Menge auch total beschränkt.
Wähle ein endliches \(\delta\)-Netz \(F\subseteq K\).
Für jedes der endlich vielen \(c\in F\) gibt es nach der
eben bewiesenen Eigenschaft ein \(U_c\in\mathcal U\) mit
\(B_K(c,\delta)\subseteq U_c\). Die Auswahl dieser endlich
vielen Mengen ist eine endliche Auswahl. Es folgt
\[
 K=\bigcup_{c\in F}B_K(c,\delta)
   \subseteq\bigcup_{c\in F}U_c\subseteq K.
\]
Also überdecken endlich viele Elemente von \(\mathcal U\)
den Raum. Damit ist jede folgenkompakte metrische Menge kompakt.

\section*{Was eine offene Überdeckung über Folgen verrät}

Für die umgekehrte Richtung sei \(K\) kompakt. Zunächst folgt
sofort totale Beschränktheit: Für jedes \(\varepsilon>0\)
bilden alle Kugeln \(B_K(c,\varepsilon)\), \(c\in K\),
eine offene Überdeckung. Die Kugeln sind offen, weil bei
\(x\in B_K(c,\varepsilon)\) jede Kugel um \(x\) mit Radius
\(\varepsilon-d(x,c)>0\) darin liegt; dies folgt aus der
Dreiecksungleichung. Eine endliche Teilüberdeckung liefert
die Zentren eines endlichen \(\varepsilon\)-Netzes.

Zur Folgenkompaktheit betrachten wir eine beliebige Folge
\((x_n)\) in \(K\). Wir suchen einen Punkt \(x\in K\),
in dessen jeder Kugel unendlich viele \emph{Indizes} der Folge
liegen. Indizes sind entscheidend: Auch eine konstante Folge
soll erfasst werden, obwohl sie nur einen einzigen Punkt trifft.

Angenommen, es gäbe keinen solchen Punkt. Dann besitzt jeder
Punkt \(x\in K\) eine Kugel, in der nur endlich viele
Folgengliederindizes liegen. Für \(V\subseteq K\) bezeichne
\(J(V)=\{n\in\mathbb N\mid x_n\in V\}\) die Menge der
betroffenen Indizes. Betrachte die Familie aller solchen Kugeln:
\[
 \mathcal W=
 \bigl\{B_K(x,r)\bigm|x\in K,\ r>0,
       \ J(B_K(x,r))\text{ ist endlich}\bigr\}.
\]
Diese Familie ist eine offene Überdeckung von \(K\), denn
jeder Punkt liegt in einer der vorausgesetzten Kugeln.
Wegen Kompaktheit reichen endlich viele davon aus. Jede dieser
Kugeln enthält jedoch nur endlich viele Indizes, und auch deren
endliche Vereinigung ist endlich. Da die Kugeln ganz \(K\)
überdecken, müssten sie zugleich alle Indizes \(\mathbb N\)
erfassen. Das ist unmöglich.

Es gibt somit einen Punkt \(x\in K\), für den jede Kugel
unendlich viele Indizes enthält. Setze
\[
 \begin{aligned}
 j_0&=\min\{n\in\mathbb N\mid d(x_n,x)<1\},\\
 j_{k+1}&=\min\{n>j_k\mid d(x_n,x)<1/(k+2)\}.
 \end{aligned}
\]
Die Mengen sind nichtleer; im zweiten Schritt bleiben aus
einer unendlichen Indexmenge nach Ausschluss des endlichen
Anfangsabschnitts noch Indizes übrig. Die Minima und die
Rekursion bestimmen eine streng wachsende Folge ohne weitere
Auswahlentscheidungen. Es gilt
\[
 d(x_{j_k},x)<\frac1{k+1}\longrightarrow0.
\]
Also konvergiert \((x_{j_k})\) gegen \(x\in K\).
Damit ist \(K\) folgenkompakt.

\section*{Drei Beschreibungen derselben Eigenschaft}

Alle Richtungen sind nun bewiesen. Für jede Teilmenge \(K\)
eines metrischen Raums sind folgende Aussagen äquivalent:
\begin{enumerate}
 \item Jede offene Überdeckung von \(K\) besitzt eine endliche
       Teilüberdeckung.
 \item Jede Folge in \(K\) besitzt eine Teilfolge, die gegen
       einen Punkt von \(K\) konvergiert.
 \item \((K,d_K)\) ist vollständig, und zu jedem
       \(\varepsilon>0\) besitzt \(K\) ein endliches
       \(\varepsilon\)-Netz.
\end{enumerate}
Dies ist \FormulaRefAuto{MetricCompactnessEquivalences}[thm].
Die nummerierten Aussagen stehen unter \CompactnessMainLink{},
die Einzelbeweise unter
\CompactnessProofsLink{}.

Die erste Beschreibung lässt aus lokalen offenen Umgebungen
endlich viele für den ganzen Raum gewinnen. Die zweite macht
aus beliebigen Folgen konvergente Teilfolgen. Die dritte
zerlegt die Eigenschaft in zwei einzeln prüfbare Bedingungen.
Je nach Aufgabe kann eine dieser Beschreibungen besonders
nützlich sein, während die anderen anschließend durch den
Äquivalenzsatz zur Verfügung stehen.

Die diskrete Menge \(\mathbb N\) und das offene Intervall
\((0,1)\) zeigen dabei zwei verschiedene Hindernisse.
Im ersten Fall lassen sich unendlich viele Punkte mit festem
Mindestabstand finden. Im zweiten Fall rücken Punkte zusammen,
aber der zugehörige Grenzwert fehlt im Raum. Totale Beschränktheit
schließt das erste Hindernis aus, Vollständigkeit das zweite.
